% 第二章 相关技术介绍
% 编译方式：xelatex -> biber -> xelatex -> xelatex

\chapter{相关技术介绍}
\label{chap:related-technologies}

在线视频分享平台的实现需要多个技术环节协同完成。系统既要处理用户、视频、评论、审核等结构化业务数据，也要支持大文件上传、视频转码、流媒体播放、全文检索、异步任务和智能内容审核等能力。因此，本章围绕论文设计与实现中的关键技术进行说明。

\section{微服务架构与服务治理技术}
\label{sec:microservice-governance}

\subsection{微服务架构思想}
\label{subsec:microservice-architecture}

在功能较少的系统中，单体架构能够降低开发和部署成本。但在线视频分享平台涉及用户认证、视频投稿、资源转码、搜索检索、互动评论、后台审核和内容安全等多个业务领域。如果所有功能都集中在一个应用中，代码耦合度会逐渐升高，后续扩展和维护也会变得困难。

微服务架构强调的是围绕业务能力来拆分服务，每个服务负责相对独立的职责，并且借助轻量级通信方式来进行协作\cite{lewis2014microservices}。对于视频平台而言，视频资源处理、普通业务管理以及智能审核任务在数据形态、处理耗时和访问压力这些方面存在明显差异，采用服务化拆分有助于降低模块之间的耦合程度。

服务拆分并不是简单按照代码目录来划分，而是要结合业务职责、数据边界以及访问压力去进行分析。表~\ref{tab:service-split-principles} 把服务拆分时需要关注的因素做了概括。

\begin{table}[htbp]
  \centering
  \caption{服务拆分原则概括}
  \label{tab:service-split-principles}
  \begin{tabular}{p{0.22\textwidth}p{0.68\textwidth}}
    \hline
    \textbf{原则} & \textbf{说明} \\
    \hline
    职责清晰 & 每个服务围绕明确业务能力设计，避免一个服务承担过多无关功能 \\
    高内聚低耦合 & 服务内部功能联系紧密，服务之间减少不必要的直接依赖 \\
    数据边界明确 & 尽量让核心数据归属于对应服务，降低跨服务频繁读写带来的复杂度 \\
    便于扩展 & 对访问压力大或处理耗时的模块，可单独扩展和部署 \\
    运维可控 & 服务数量需要与项目规模匹配，避免过度拆分导致管理成本过高 \\
    \hline
  \end{tabular}
\end{table}

\subsection{服务构建与运行基础}
\label{subsec:spring-boot-cloud}

Spring Boot 是 Java Web 开发当中比较常用的一种服务构建框架。它借助自动配置、起步依赖以及内嵌服务器这些机制来减少重复配置，让开发者能够较快地搭建起独立运行的后端服务\cite{springbootdocs}。在这类系统当中，用户服务、资源服务、互动服务以及后台管理服务基于 Spring Boot 来进行构建。

当系统拆分出多个服务之后，还需要去处理服务注册、配置管理、网关路由以及远程调用这些问题。Spring Cloud 提供了一套微服务开发组件，用来支持分布式系统当中的服务治理能力\cite{springclouddocs}。所以说，Spring Boot 更侧重于单个服务的快速构建，而 Spring Cloud 则更侧重于多个服务之间的协同管理。

\subsection{服务注册、网关与远程调用}
\label{subsec:registry-gateway-feign}

在微服务环境当中，服务实例可能会因为重启、扩容或者部署迁移而改变地址。要是调用方直接写死服务地址，那就会降低系统的灵活性和可维护性。服务注册与发现机制能让服务启动之后向注册中心登记自身信息，调用方再借助服务名去获取可用的实例。Nacos 支持服务发现和配置管理，比较适合用在微服务系统里的注册中心和配置中心\cite{nacosdocs}。

系统对外访问通常还需要一个统一的入口。Spring Cloud Gateway 是 Spring Cloud 体系当中的网关组件，它依据路由规则把请求转发到不同的服务，并且在网关层统一去处理跨域、鉴权、日志以及内部接口隔离这些通用逻辑\cite{springgatewaydocs}。引入网关减少客户端直接感知多个内部服务所带来的复杂度。

服务之间的同步调用借助 OpenFeign 来实现。OpenFeign 是一个声明式的 HTTP 客户端，开发者通过接口去描述远程调用关系，由框架完成请求的发送以及响应的转换\cite{springopenfeign}。不过同步调用比较适合那些耗时较短、需要马上返回结果的场景，而视频转码和 AI 审核这类长耗时任务则更适合使用异步消息机制。

在跨服务写操作的场景当中，单个服务的本地事务很难覆盖完整的业务链路。Seata 是一个面向微服务架构的分布式事务解决方案，它能通过全局事务协调机制去处理跨服务的数据提交与回滚问题\cite{seatadocs}。本系统在公共基础模块当中引入了 Seata，并在涉及投稿审核、互动计数以及用户资产变更这类关键写操作时，为数据一致性控制提供技术方面的支撑。

表~\ref{tab:microservice-tech-summary} 把微服务治理相关的技术做了一个概括总结。

\begin{table}[htbp]
  \centering
  \caption{微服务治理相关技术作用}
  \label{tab:microservice-tech-summary}
  \begin{tabular}{p{0.24\textwidth}p{0.32\textwidth}p{0.34\textwidth}}
    \hline
    \textbf{技术} & \textbf{主要作用} & \textbf{适用环节} \\
    \hline
    Spring Boot & 快速构建独立运行的后端服务 & 各业务服务基础框架 \\
    Spring Cloud & 提供微服务治理组件 & 多服务协同与治理 \\
    Nacos & 服务注册发现、配置管理 & 服务地址管理与配置集中维护 \\
    Spring Cloud Gateway & 统一系统入口与请求转发 & 前端访问、鉴权、内部接口保护 \\
    OpenFeign & 声明式远程调用 & 服务之间的短耗时同步调用 \\
        Seata & 分布式事务协调、全局提交与回滚控制 & 跨服务写操作的数据一致性控制 \\
    \hline
  \end{tabular}
\end{table}

\section{数据存储、缓存与检索技术}
\label{sec:data-storage-cache-search}

\subsection{关系型数据库与数据访问层}
\label{subsec:mysql-mybatis}

关系型数据库比较适合保存结构清晰、关系明确的业务数据。MySQL 是一种常用的关系型数据库管理系统，它拥有成熟的事务、索引以及查询能力\cite{mysqldocs}。在在线视频分享平台当中，用户信息、视频基本信息、分类信息、评论记录以及审核记录这些数据都属于典型的结构化业务数据，适合由关系型数据库来统一管理。

为了避免在业务代码当中直接拼接大量 SQL，Java 后端通常会引入数据持久层框架。MyBatis 支持通过 XML 或者注解来编写 SQL，并且把查询结果映射为 Java 对象\cite{mybatisdocs}。跟完全自动化的 ORM 框架相比，MyBatis 保留了较强的 SQL 控制能力，比较适合查询条件比较灵活的业务系统。

\subsection{Redis 缓存与临时状态管理}
\label{subsec:redis-cache}

缓存技术主要用来提升热点数据的访问速度，同时减少数据库的压力。Redis 是一种基于内存的数据存储系统，它支持字符串、列表、集合、有序集合以及哈希这些数据结构\cite{redisdocs}。由于其读写速度比较快、数据结构也比较灵活，所以经常被用在验证码、登录状态、访问计数、分类缓存以及临时任务状态等场景当中。

需要留意的是，Redis 更适合用来保存热点数据和临时状态，不能简单地替代关系型数据库。核心业务数据仍保存在 MySQL 当中，而缓存层则用来提升访问效率或者辅助异步流程。让正式数据和临时数据各自承担起适宜的职责。

\subsection{Elasticsearch 全文检索}
\label{subsec:elasticsearch-search}

视频平台的搜索功能不仅需要根据标题和简介进行关键词匹配，还要考虑分词、相关性排序以及搜索结果召回这些问题。普通数据库虽然也能做模糊查询，但在文本内容比较多、检索条件比较复杂的时候，效率和效果都会受到限制。

Elasticsearch 是一个分布式的搜索与分析引擎，它基于倒排索引来实现全文检索，经常被用在搜索、日志分析以及复杂查询这些场景当中\cite{elasticsearchdocs}。在内容类的平台里，Elasticsearch 通常作为业务数据库之外的检索索引层来使用。MySQL 负责保存权威的业务数据，Elasticsearch 则负责提升搜索体验。

\begin{table}[htbp]
  \centering
  \caption{数据处理相关技术对比}
  \label{tab:data-tech-summary}
  \begin{tabular}{p{0.20\textwidth}p{0.30\textwidth}p{0.38\textwidth}}
    \hline
    \textbf{技术} & \textbf{主要特点} & \textbf{适用场景} \\
    \hline
    MySQL & 结构化存储、事务能力成熟 & 用户、视频、评论、审核等核心业务数据 \\
    MyBatis & SQL 可控、便于对象映射 & 数据访问层封装和复杂查询组织 \\
    Redis & 内存读写快、数据结构灵活 & 验证码、登录态、缓存、计数和临时状态 \\
    Elasticsearch & 倒排索引、全文检索能力强 & 视频标题、简介、标签等内容检索 \\
    \hline
  \end{tabular}
\end{table}

\section{视频上传、转码与播放技术}
\label{sec:video-upload-transcode-play}

\subsection{大文件分片上传}
\label{subsec:chunk-upload}

视频文件通常体积比较大，要是一次性上传完整的文件，很容易受到网络波动、接口超时以及浏览器限制的影响。分片上传的基本思路是把一个大文件划分成多个较小的片段，由客户端逐个进行上传，服务端记录上传进度，等所有分片都到达之后再完成合并。

分片上传能够降低单次请求失败所带来的重传成本，也便于系统去记录上传状态。对于视频投稿这个场景来说，分片上传能够提升大文件上传的稳定性，这也是后续转码和审核流程能够顺利执行的基础。

\subsection{FFmpeg 媒体转码}
\label{subsec:ffmpeg-transcode}

用户上传的视频可能来自不同的设备和软件，视频格式、编码方式、分辨率以及码率并不统一。要是系统直接提供原始文件进行播放，就有可能出现浏览器兼容性差、加载慢或者播放失败这些问题。因此，视频平台通常需要对上传的文件进行统一的转码。

FFmpeg 是一个常用的开源音视频处理工具，它能完成视频转码、格式转换、音频提取、截图以及切片这些操作\cite{ffmpegdocs}。借助转码把不同来源的视频转换成适合在线播放的格式，并且根据需要去生成封面、音频分析文件或者流媒体播放资源。

媒体转码属于比较耗时的操作，不适合放在用户上传请求当中同步完成。更合理的方式是，在用户提交视频之后生成后台任务，由异步消费者去完成转码处理，然后再把处理状态回写到数据库或者缓存当中。

\subsection{HLS 流媒体播放}
\label{subsec:hls-streaming}

在线播放视频的时候，要是直接加载完整的视频文件，会增加初始等待时间，也不利于弱网环境下的播放体验。HLS 是一种基于 HTTP 的流媒体传输协议，它把视频组织成播放列表和多个媒体分片，播放器先读取播放列表，再根据需要去加载对应的片段\cite{rfc8216hls}。

HLS 常见的文件包括 \texttt{.m3u8} 播放列表和 \texttt{.ts} 媒体分片。用户看到的是连续播放的视频，而系统实际提供的是一组按需加载的片段资源。这种方式有利于降低视频初始加载的压力，也便于后续扩展多清晰度的播放。

\subsection{浏览器视频播放能力}
\label{subsec:html5-video}

现代浏览器借助 HTML5 的 \texttt{video} 元素来提供视频播放能力，支持播放、暂停、音量控制、进度控制以及全屏这些基本操作\cite{whatwgvideo}。不过，完整的视频播放体验不仅依赖浏览器的控件，还需要服务端提供兼容的视频格式、稳定的资源路径以及合理的加载策略。

视频播放链路能概括为“上传文件、进行转码处理、生成播放资源、由浏览器加载并播放”。各个技术环节的作用如表~\ref{tab:video-tech-summary-revised} 所示。

\begin{table}[htbp]
  \centering
  \caption{视频处理相关技术作用}
  \label{tab:video-tech-summary-revised}
  \begin{tabular}{p{0.24\textwidth}p{0.36\textwidth}p{0.30\textwidth}}
    \hline
    \textbf{技术环节} & \textbf{解决的问题} & \textbf{作用} \\
    \hline
    分片上传 & 大文件上传失败后重传成本高 & 提升上传稳定性 \\
    FFmpeg 转码 & 原始格式不统一、播放兼容性差 & 生成统一播放资源 \\
    HLS 分片 & 完整加载视频等待时间长 & 支持按需加载和在线播放 \\
    HTML5 video & 浏览器页面播放控制 & 提供前端播放能力 \\
    \hline
  \end{tabular}
\end{table}

\section{前端页面开发与接口调用技术}
\label{sec:frontend-tech}

系统前端包含普通用户端和后台管理端两个部分，页面开发采用 Vue 3 与 Vite 作为基础工具链。Vue 3 负责组件化页面构建，Vite 用于提升开发启动和构建效率；Vue Router 用于管理首页、播放页、创作中心、个人中心和后台管理页面之间的路由跳转；Pinia 用于维护用户登录态、页面共享状态和部分业务状态\cite{vueguide,viteguide,vuerouterdocs,piniadocs}。

后台管理端需要较多表格、表单、弹窗、分页和审核操作组件，因此采用 Element Plus 提供基础 UI 组件；前端与后端服务之间通过 Axios 发起 HTTP 请求，并由网关统一转发到对应微服务\cite{elementplusdocs,axiosdocs}。这种前后端分离方式能够降低页面展示逻辑与后端业务逻辑之间的耦合，也便于普通用户端和管理端分别维护。

\section{异步通信与任务处理技术}
\label{sec:async-message-task}

\subsection{异步处理机制}
\label{subsec:async-mechanism}

在 Web 系统当中，并不是所有的操作都适合马上完成并返回结果。对于视频转码、文件清理、AI 审核这类比较耗时的任务，要是让前台请求一直等待，很容易造成接口超时，也会降低用户体验。数据密集型应用设计当中强调，可靠性与可维护性需要结合任务状态、故障恢复以及系统解耦来加以考虑\cite{kleppmann2017ddia}。

异步处理的基本思想是：用户先提交任务，系统记录任务状态，然后由后台消费者执行具体任务。用户或管理员能通过状态查看处理进度。这样既能减少前台请求等待时间，也能提高复杂任务处理的稳定性。

\subsection{RabbitMQ 消息队列}
\label{subsec:rabbitmq}

RabbitMQ 是一种消息队列中间件，支持交换机、队列、路由键、消息确认、持久化和死信队列等机制\cite{rabbitmqdocs}。在视频平台中，RabbitMQ 负责连接“任务生产者”和“任务消费者”，例如上传服务产生转码任务，资源服务或转码消费者异步处理视频文件。

消息队列能够降低服务之间的直接耦合程度，但同时也会带来消息重复、消费失败、任务幂等以及状态补偿这些问题。因此，系统在使用消息队列的时候，不能只是发送消息，还需要配合任务状态表、失败重试以及异常记录这些机制。

\subsection{任务状态与一致性控制}
\label{subsec:task-status-consistency}

异步任务跟同步请求的主要区别在于，异步任务通常要经历待处理、处理中、处理成功、处理失败等多个状态。微服务模式当中也强调，事件驱动的流程往往需要配合状态记录、重试以及补偿机制来使用\cite{richardson2018microservices}。

在视频平台当中，转码任务、审核任务以及文件清理任务都需要有明确的状态流转。任务状态既能够在前端展示处理进度，也能在后台排查异常以及避免重复处理。对于那些不要求强一致性、不需要马上完成的业务也能采用最终一致和异步补偿的策略，以此来减少强事务带来的复杂度。

\begin{table}[htbp]
  \centering
  \caption{同步调用与异步通信对比}
  \label{tab:sync-async-summary-revised}
  \begin{tabular}{p{0.20\textwidth}p{0.34\textwidth}p{0.34\textwidth}}
    \hline
    \textbf{比较项} & \textbf{同步调用} & \textbf{异步通信} \\
    \hline
    适用任务 & 查询、登录、短时间业务处理 & 转码、审核、文件处理等耗时任务 \\
    用户体验 & 需要等待处理结果 & 前台快速返回，后台继续执行 \\
    服务耦合 & 调用方依赖被调用方即时响应 & 生产者和消费者相对解耦 \\
    异常处理 & 失败可立即返回 & 需要记录状态、重试和补偿 \\
    \hline
  \end{tabular}
\end{table}

\section{智能内容审核技术}
\label{sec:intelligent-content-audit-revised}

\subsection{视频内容审核思路}
\label{subsec:audit-basic-idea}

内容审核的目标是在内容正式发布之前识别出可能存在的违规或者风险信息。传统审核主要依赖人工判断或者规则匹配，而算法内容审核通常需要跟人工判断、规则体系以及申诉机制结合起来使用\cite{gorwa2020algorithmic}。

视频内容比普通的文本和图片要更复杂。它同时包含了画面、语音、背景声音以及上下文语义，某些风险信息可能隐藏在语音内容或者关键画面当中。智能内容审核的作用就是借助语音识别、声音事件识别以及多模态理解这些技术，为人工审核提供辅助依据，而不是完全替代人工判断。

\subsection{语音识别与声音事件识别}
\label{subsec:speech-audio-event}

视频当中的语音内容往往带有比较重要的语义信息。Whisper 是一种基于大规模弱监督训练的语音识别模型，它在多语言语音识别任务上有着较好的表现\cite{radford2023whisper}。借助语音识别，把视频当中的人声转换成文本，再交由规则或者模型做进一步分析。

除了语音内容之外，声音事件也可能提示视频存在风险。比如说尖叫、爆炸声、警报声这类异常声音，可能对应着特殊的场景。AudioSet 是一个大规模音频事件数据集，它为声音事件识别方面的研究提供了基础\cite{gemmeke2017audioset}。YAMNet 这类声音事件分类模型关注的是“发生了什么声音”，这与语音识别关注“说了什么内容”有所不同。

\subsection{多模态理解与人机协同审核}
\label{subsec:multimodal-human-machine}

多模态语义理解指的是模型同时去处理图像、文本、语音等多种信息，并且综合判断内容的含义。Qwen-VL 是一个视觉语言模型，它能够结合图像和文本进行理解，用于图像描述、视觉问答、文字识别以及多模态推理等任务\cite{bai2023qwen}。在视频审核场景，从视频里抽取出关键帧，再结合语音转写出来的文本进行综合分析。

考虑到智能模型可能存在误判、漏判或者推理结果不稳定这些问题，内容审核系统更适合采用“AI 做初审，人工做终审”的方式。AI 审核负责快速发现风险片段、生成风险等级以及审核建议；管理员再根据模型结果、视频内容以及平台规则来做出最终的决策。这种方式既能提高审核效率，也能保留人工判断的灵活性。

\begin{table}[htbp]
  \centering
  \caption{智能审核相关技术作用}
  \label{tab:audit-tech-summary-revised}
  \begin{tabular}{p{0.24\textwidth}p{0.34\textwidth}p{0.32\textwidth}}
    \hline
    \textbf{技术} & \textbf{分析对象} & \textbf{审核作用} \\
    \hline
    Whisper & 视频语音内容 & 将语音转换为文本，辅助识别风险表达 \\
    YAMNet & 背景声音和事件声音 & 发现异常声音事件，提供辅助线索 \\
    Qwen-VL & 关键帧与文本信息 & 综合画面和语义，生成风险判断依据 \\
    人工审核 & 视频整体内容与平台规则 & 对 AI 结果进行复核并作出最终决策 \\
    \hline
  \end{tabular}
\end{table}